\documentclass[10pt,a4paper]{article}
\usepackage{cv-style}
\geometry{margin=0.62in}
\hypersetup{
  pdftitle={CV - Filippo Lollato - Tirocinio AI FERCAM},
  pdfauthor={Filippo Lollato}
}
\setlength{\parindent}{0pt}
\setlength{\parskip}{0.3ex}
\titlespacing*{\section}{0pt}{1.05ex}{0.55ex}

\begin{document}

\name{Filippo Lollato}
\contact{\href{mailto:contact@lollato.dev}{contact@lollato.dev} \enskip\textbullet\enskip \href{https://lollato.dev}{lollato.dev} \enskip\textbullet\enskip \href{https://github.com/xupremix}{github.com/xupremix} \enskip\textbullet\enskip Trento, Italia}

\section{Profilo}
Studente magistrale in Artificial Intelligence Systems e laureato in Informatica. Mi interessa applicare l'AI a problemi operativi concreti, integrando modelli e agenti nei flussi di lavoro e nei sistemi software esistenti. Nei progetti universitari ho lavorato su integrazione LLM, sistemi multi-agente e pipeline ML; cerco un tirocinio in cui contribuire alla sperimentazione e all'adozione pratica di soluzioni AI.

\section{Istruzione}
\noindent\textbf{Laurea Magistrale in Artificial Intelligence Systems} \hfill \textsl{2025--26 a oggi}\\
\noindent Università di Trento \hfill Trento, Italia

\vspace{0.45ex}
\noindent\textbf{Laurea Triennale in Informatica} \hfill \textsl{2022--23 -- 2024--25}\\
\noindent Università di Trento \hfill Trento, Italia\\
\noindent Tesi: \textit{incin}, framework di machine learning in Rust con controlli di forma e invarianti dei tensori a tempo di compilazione.

\section{Esperienza}
\noindent\textbf{Tutor universitario / Assistente alla didattica} \hfill \textsl{2024--25 a oggi}\\
\noindent Università di Trento \hfill Trento, Italia
\begin{itemize}
  \item Tutoraggio in Rust, Programmazione Avanzata e Compilatori dal terzo anno della laurea triennale.
  \item Supporto agli studenti nel comprendere requisiti, interpretare errori e risolvere problemi di programmazione; esperienza nel comunicare concetti tecnici in modo chiaro.
\end{itemize}

\section{Progetti rilevanti}
\noindent\textbf{ASA Autobots -- Sistema multi-agente per consegne} \hfill \textsl{Node.js, BDI, PDDL, LLM}\\
\begin{itemize}
  \item Co-responsabile di un team di due persone per un workflow di consegna che combina agenti BDI, un coordinatore basato su LLM e un planner PDDL.
  \item Contributo all'architettura del sistema e all'integrazione dei componenti agentici e delle interfacce tra servizi.
\end{itemize}

\noindent\textbf{incin -- Framework ML con controlli a tempo di compilazione} \hfill \textsl{Rust, Machine Learning}\\
\begin{itemize}
  \item Progetto di tesi personale: esplora controlli statici per forme e invarianti dei tensori, per rendere il codice ML più facile da verificare e ragionare.
  \item Obiettivo di lungo periodo: fornire regole architetturali e feedback del compilatore che aiutino agenti AI di ricerca a generare codice ML corretto.
\end{itemize}

\noindent\textbf{Classificazione di intenti e slot} \hfill \textsl{Python, PyTorch, BERT, Hugging Face}\\
\begin{itemize}
  \item Realizzata una pipeline di addestramento e valutazione su ATIS; il modello congiunto ha ottenuto un'accuratezza di 0,9765 nella classificazione degli intenti.
\end{itemize}

\section{Competenze}
\cvskill{AI e integrazione}{Integrazione LLM e AI nei workflow, sistemi multi-agente, BDI, PDDL, tool/function calling, API REST}
\cvskill{Sviluppo e strumenti}{Python, Node.js, TypeScript, Rust, C++20, PyTorch, Hugging Face, Git, Linux, Docker, CI/CD}
\cvskill{Lingue}{Italiano madrelingua; inglese Cambridge B2 First, Grade A (livello QCER C1)}

\end{document}
